% Compile beside math-review.sty; replace all visible insertion text for publication.
% Rename and regroup functional headings around the actual mathematical thread.
% Shared files do not determine section hierarchy; no environment/count quota applies.
\documentclass[12pt,reqno]{amsart}
\usepackage{math-review}
\title[Part V: Thematic Review (concise)]{[Review topic]\\
Part V: A Thematic Review (concise)}
\author{[Author name]}
\date{[Prepared date; literature cutoff date]}
\begin{document}
\begin{abstract}
\placeholder{Describe the central problem and perspective of a short thematic review, selecting the main answers, the ideas behind them and the Problems needed for an overall understanding.}
\end{abstract}
\maketitle
\tableofcontents

\section{Introduction}\label{sec:introduction}
\placeholder{State the concrete central problem, its importance and this review's mathematical perspective. Identify the larger thematic relationship that connects the selected results to the remaining questions. Functional labels such as background and methods do not by themselves provide that relationship.}
\placeholder{Explain why the language in Section~\ref{sec:setting} is enough for the selected answers in Section~\ref{sec:results}; how Section~\ref{sec:ideas} explains what makes those answers possible; and why their precise reach leads to Section~\ref{sec:problems}. Use the actual mathematical relationships and update the references when regrouping sections.}

\section{Setting and notation}\label{sec:setting}
\placeholder{Define only the objects and notation needed for the chosen main results and Problems, with the relevant conventions. Select this language from the review's own thread, not from the full foundations or model inventory of another part.}

\section{Main results and their relations}\label{sec:results}
\placeholder{Select the principal answers to the central problem and organize them by their mathematical relationship. Combine local variants, remove secondary branches and repeated scope reminders, and retain every necessary condition of the statements that remain.}
\begin{theorem}[A principal answer]\label{thm:principal}
\placeholder{State a verified main conclusion with all assumptions, quantifiers, range and qualifications, together with its source and controlling version. Reuse a shared complete statement's canonical body.}
\end{theorem}
\placeholder{Explain what Theorem~\ref{thm:principal} answers and how it relates to the other selected conclusions. A precise comparison and citation may replace a technical detour; do not delegate indispensable assumptions to another volume.}

\section{Ideas behind the results}\label{sec:ideas}
\placeholder{For the selected main results, explain the main idea, what it transforms or resolves, and how it supports the conclusion. A compact connected paragraph is usually enough. Name precise external inputs and their role where necessary; method names alone are not an explanation.}
\placeholder{Retain a short derivation only if it significantly clarifies this thread. Do not import the complete proof families, technical lemma inventory or model calculations of Part III or I. Results and their method ideas may instead be interwoven in the preceding section.}

\section{Problems and further directions}\label{sec:problems}
\placeholder{Choose the core unresolved targets serving this review's own thread; do not automatically copy an entire frontier volume or imply omitted topics are settled.}
\begin{problem}[A direction arising from the main results]\label{prob:core}
\placeholder{State a precise sourced or explicitly review-proposed target with its objects, assumptions and quantifiers. Keep history and lengthy progress discussion outside this statement.}
\end{problem}
\placeholder{After Problem~\ref{prob:core}, state its formulation source and cutoff status, why it matters, the decisive verified progress and the exact remaining scope or verification limits. Preserve conditions and status wherever the Problem is shared. Do not infer openness from search failure.}
\placeholder{For a narrowly scoped topic whose settlement is verified, replace the open-Problem scaffold with the scope and settlement evidence rather than inventing a frontier. Close by relating the selected results, method ideas and directions to the original question.}

\templatereferences
\end{document}
